\scanpage{11}
% L01–03
\uline{Prop} \[
\psi_n^{-1}\bigl(\underbrace{\mathbb E^*\times\mathbb E^{n-1}}_{\substack{(a_0,\text{--},a_{n-1})\\\text{tels que }a_0\text{ inv.}}}\bigr)=Q_n^0
\]
% L04
et
% L05
\diagram{\begin{tikzcd}[ampersand replacement=\&]
 Q_n^0 \arrow[r,"\sim"] \& \mathbb E^*\times\mathbb E^{n-1}
\end{tikzcd}}

% L06–09
\deleted{En effet,} [Il est immédiat que $Q_n^0\simeq$ ``ss-groupes finis étales de $G_a$ de rg $p^n$'', ce qui donne une interprétation géom de ce foncteur \unsure{ouvert|couvert|constant} $\mathbb E^*\times\mathbb E^{n-1}$, en termes de noyaux d'hom $f_a$ \ldots]

% L10–11
En effet, on a un \deleted{foncteur} hom en sens inverse
% L11–12
\diagram{\begin{tikzcd}[ampersand replacement=\&]
Q_n^0 \& E^*\times\mathbb E^{n-1}\subset\mathbb E^n \arrow[l]\\
f_a \& a \arrow[l,mapsto] \& \ldots
\end{tikzcd}}

% L13–14
\uline{Thm} $\unsure{\mathbb Q|Q}_n\longrightarrow\mathbb E^n$ est un isomorphisme, en d'autres termes
% L15
\[ \mathbb F_p[X_1,\text{--},X_n]^{\Gamma_n}\simeq\mathbb F_p[A_0,\text{--},A_{n-1}] \]

% L16–19
\unsure{On|Il|Dès} sait déjà que $Q_n\mathbin{-}\mathbb E_{\deleted{\unsure{n}}}^{n}$ est birationnel, il suffit de prouver qu'il est fini \inserted{ou même propre} -- comme $Q_n$ est intègre et $E_n$ normal, on conclut que c'est un iso.
% L19–21
\deleted{Il suffit de prouver} \inserted{Ou aussi il suffit} \deleted{que $\mathbb E^n\xrightarrow{A_*}\mathbb E^n$} de prouver que $Q_n\to\mathbb E^n$ \unsure{est|$n$} est quasi-fini (car \unsure{sera|ser.} une imm. ouverte par le
% archive sheet number: 10
